\section{Introduction}
\label{sec:intro}

The landscape of cybersecurity is undergoing a fundamental transformation driven by the emergence of Large Language Models (LLMs) capable of autonomous security reasoning. Unlike traditional vulnerability scanners that rely on static signature matching, LLM-powered "reasoning agents" can simulate the tactical decision-making of human penetration testers—interpreting complex tool outputs, adapting to defensive responses, and chaining minor misconfigurations into multi-stage attack paths. As organizations move toward continuous, automated security assessments, the reliability of these agents has become a paramount concern.

\subsection{The Retrieval Gap in Autonomous VAPT}
Despite the promise of autonomous red teaming, a critical bottleneck has emerged: the \textbf{Retrieval Gap}. Standard Retrieval-Augmented Generation (RAG) systems, which form the "knowledge backbone" of these agents, typically rely on keyword-based or naive vector-based search. While effective for simple question-answering, these methods fail in the nuanced domain of vulnerability intelligence. Specifically, they struggle to:
\begin{enumerate}
    \item \textbf{Identify Structural Links}: Many vulnerabilities are functionally related (e.g., helper library flaws or transitive dependencies) but semantically distinct, leading to missed attack vectors in flat vector stores.
    \item \textbf{Prioritize Operational Actionability}: Semantic similarity does not equal operational priority. A CVE with a high similarity score but no public exploit is often less relevant than a slightly less similar CVE that is actively being exploited in the wild.
    \item \textbf{Overcome Query Ambiguity}: Security queries are often imprecise (e.g., "memory corruption in apache"). Standard RAG systems provide a single-shot response, lacking the agentic persistence to reformulate queries when initial results are insufficient.
\end{enumerate}

\subsection{Our Contribution: VulnRAG}
In this paper, we introduce \textbf{VulnRAG}, an agentic RAG framework designed to bridge the Retrieval Gap through structural discovery and active self-correction. Unlike generic reflective RAG systems (e.g., Self-RAG), VulnRAG’s agentic loop is specifically architected for security-domain challenges, incorporating automated Common Platform Enumeration (CPE) scope expansion and vulnerability lifecycle metadata into the feedback loop. Our core contributions are:

\begin{itemize}
    \item \textbf{Structural-Semantic Hybrid Retrieval}: We implement a novel retrieval pipeline that integrates dense embeddings with structural traversal of a Directed Acyclic Graph (DAG) derived from CVE relationship metadata. This enables the discovery of functionally linked vulnerabilities that traditional RAG misses.
    \item \textbf{Multi-Factor Semantic Reranking (MFSR)}: We propose a security-aware relevance function that integrates Cross-Encoder semantic scores with CVSS severity, exploit availability, and disclosure recency, ensuring the most actionable intelligence is prioritized in the LLM context.
    \item \textbf{Zero-Shot Agentic Self-Correction}: We design a training-free feedback loop where a decoupled Evaluator agent autonomously optimizes retrieval parameters and query expansions, matching the performance of fine-tuned reflective models in a provider-agnostic manner.
    \item \textbf{Complexity-Aware Multi-Tier Routing}: We demonstrate a tiered execution model that optimizes the performance-to-cost ratio, achieving flagship-level reasoning quality at near-Flash tier costs through intelligent task escalation.
\end{itemize}

The remainder of this paper is structured as follows: Section~\ref{sec:background} reviews related work in RAG and autonomous security; Section~\ref{sec:threat} formalizes our threat model; Section~\ref{sec:design} and \ref{sec:optimization} detail the VulnRAG architecture and agentic loops; Section~\ref{sec:reranking} formalizes the MFSR algorithm; Section~\ref{sec:evaluation} presents our empirical results; and Section~\ref{sec:discussion} discusses ethical considerations and limitations.
